\documentclass[11pt]{article}
\usepackage[a4paper,margin=2.5cm]{geometry}
\usepackage{amsmath}
\usepackage{amssymb}
\usepackage{enumitem}
\usepackage{parskip}
\setlist[enumerate,2]{label=(\alph*), itemsep=1pt, topsep=2pt}
\setlength{\parindent}{0pt}
\newcommand{\correct}[1]{\textbf{#1}\;$\checkmark$}
\begin{document}
\section*{Auctions}
\begin{enumerate}[itemsep=6pt]
  \item In a second-price (Vickrey) auction, the winner ...
  \begin{enumerate}
    \item pays the average of all bids
    \item is the highest bidder and pays the second-highest bid
    \item is chosen at random
    \item is the highest bidder and pays their own bid
  \end{enumerate}
  \item In a second-price auction, bidding your true value is ...
  \begin{enumerate}
    \item the same as bidding zero
    \item a weakly dominant strategy
    \item never optimal
    \item only optimal if you know the others' values
  \end{enumerate}
  \item A Bayesian Nash equilibrium is a strategy profile in which ...
  \begin{enumerate}
    \item the seller chooses the winner
    \item all players bid the same amount
    \item each player maximises expected utility given their type and correct beliefs about others
    \item every player knows everyone else's value
  \end{enumerate}
  \item In a first-price auction, a rational bidder usually ...
  \begin{enumerate}
    \item shades their bid below their true value
    \item always bids zero
    \item bids their true value
    \item bids above their true value
  \end{enumerate}
  \item The revenue equivalence theorem says that ...
  \begin{enumerate}
    \item first-price and second-price auctions raise the same expected revenue (under its conditions)
    \item the winner always pays the same as the loser
    \item the seller always prefers a first-price auction
    \item revenue is always zero
  \end{enumerate}
  \item Because truthful bidding is dominant, the second-price auction is ...
  \begin{enumerate}
    \item always more profitable for the seller
    \item impossible to run with sealed bids
    \item only suitable for two bidders
    \item efficient: the highest-value bidder wins
  \end{enumerate}
  \item An auction game is specified by ...
  \begin{enumerate}
    \item a single payoff matrix
    \item a discount factor
    \item the bidders' value distributions, an allocation rule and a payment rule
    \item a list of stable matchings
  \end{enumerate}
  \item In a first-price auction, the winner pays ...
  \begin{enumerate}
    \item nothing
    \item the average bid
    \item their own bid
    \item the second-highest bid
  \end{enumerate}
  \item A bidder's type in an auction is ...
  \begin{enumerate}
    \item their private value for the item
    \item their bid
    \item the reserve price
    \item the number of other bidders
  \end{enumerate}
  \item In a first-price auction with \(N\) bidders whose values are uniform on \([0,1]\), the equilibrium bid is ...
  \begin{enumerate}
    \item \(v/2\) for all \(N\)
    \item \(\frac{N-1}{N} v\)
    \item \(v\)
    \item \(0\)
  \end{enumerate}
  \item A bidder's utility in an auction is ...
  \begin{enumerate}
    \item always zero
    \item \(v_i q_i - p_i\): value times allocation, minus payment
    \item their bid
    \item the seller's revenue
  \end{enumerate}
  \item Revenue equivalence requires that the auction ...
  \begin{enumerate}
    \item always gives the item to the highest-value bidder and gives a zero-value bidder zero surplus
    \item has exactly two bidders
    \item charges a reserve price
    \item uses sealed bids
  \end{enumerate}
  \item In a second-price auction the winner's surplus is ...
  \begin{enumerate}
    \item the second-highest value
    \item zero
    \item their value minus their own bid
    \item their value minus the second-highest value
  \end{enumerate}
  \item Truthful bidding being a dominant strategy means that in a second-price auction ...
  \begin{enumerate}
    \item impossible to run sealed-bid
    \item bidders need not guess others' values to bid optimally
    \item only suitable for one bidder
    \item more profitable than any other auction
  \end{enumerate}
  \item The second-price auction is efficient because ...
  \begin{enumerate}
    \item the seller earns the most
    \item no one ever wins
    \item everyone pays the same
    \item the bidder who values the item most wins it
  \end{enumerate}
  \item For \(N\) values drawn independently and uniformly on \([0,1]\), the expected maximum is ...
  \begin{enumerate}
    \item \(\frac{1}{N}\)
    \item \(\frac{1}{2}\)
    \item \(\frac{N-1}{N+1}\)
    \item \(\frac{N}{N+1}\)
  \end{enumerate}
  \item In a first-price auction with \(N\) uniform bidders bidding \(\frac{N-1}{N} v\), the seller's expected revenue is ...
  \begin{enumerate}
    \item \(\frac{N-1}{N+1}\)
    \item \(\frac{N}{N+1}\)
    \item \(\frac{1}{N}\)
    \item \(\frac{1}{2}\)
  \end{enumerate}
  \item In a second-price auction with \(N\) uniform bidders, the seller's expected revenue equals ...
  \begin{enumerate}
    \item the expected second-highest value, \(\frac{N-1}{N+1}\)
    \item the expected highest value, \(\frac{N}{N+1}\)
    \item \(\frac{1}{2}\)
    \item zero
  \end{enumerate}
  \item A bidding strategy in an auction game is a map \(b_i: V_i \to B_i\) that sends ...
  \begin{enumerate}
    \item a bid to a payment
    \item a valuation to a payment
    \item a bidder's valuation to a bid
    \item a bid to an allocation
  \end{enumerate}
  \item The allocation rule \(q\) of an auction maps the profile of bids to ...
  \begin{enumerate}
    \item the probability with which each bidder receives the good
    \item the amount each bidder pays
    \item each bidder's private valuation
    \item the seller's reserve price
  \end{enumerate}
\end{enumerate}
\end{document}
